\documentclass[aspectratio=169,10pt]{beamer}
\usepackage{plan4ari_beamer}

\usepackage{iftex}
\ifPDFTeX%
  \usepackage[T1]{fontenc}
  \usepackage[utf8]{inputenc}
  % Times-compatible fallback for pdfLaTeX.
  \usepackage{newtxtext,newtxmath}
\else
  \usepackage{fontspec}
  % Prefer one of the reference fonts named in the official template.
  \IfFontExistsTF{Times New Roman}{%
    \setmainfont{Times New Roman}
  }{%
    \IfFontExistsTF{Nimbus Roman No9 L}{%
      \setmainfont{Nimbus Roman No9 L}
    }{%
      \IfFontExistsTF{Nimbus Roman}{%
        \setmainfont{Nimbus Roman}
      }{%
        % Metric-compatible Times-style fallback; verify against the official template before submission.
        \setmainfont{TeX Gyre Termes}
      }%
    }%
  }
\fi
\usepackage{lmodern}
\usepackage{graphicx}
\usepackage{microtype}
\usepackage{ragged2e}

\usepackage{amsmath,amssymb}
\usepackage{tikz}
\usetikzlibrary{arrows,positioning,shapes.geometric,fit,calc,matrix,shadows.blur}
\usepackage{pgfgantt}
\usepackage[table]{xcolor}
\makeatletter
\def\input@path{{./styles/}}
\makeatother
\usepackage{plan4ari_theme}
\hypersetup{colorlinks=true,linkcolor=planblue,urlcolor=planblue,citecolor=planblue}

\usepackage[
    backend=biber,
    style=verbose,
    autocite=footnote,
    doi=true,
    url=false
]{biblatex}


\usepackage[marginal]{footmisc}


% ---------- Tables, figures, lists ----------
\usepackage{booktabs,tabularx,array,multirow}
\usepackage{longtable}
\usepackage{makecell}
\usepackage{subcaption}
\usepackage{caption}
\captionsetup{labelfont=bf}


% ---------- Header/footer and page-count warning ----------
\usepackage{empheq}
\usepackage[most]{tcolorbox}
\tcbuselibrary{breakable}

\usepackage[normalem]{ulem}

\usepackage{xparse}

\usepackage{xltabular}
\keepXColumns

% Compatibility for longtable v4.24 and older: its output routine used
% infinitely shrinkable \vss glue, which LuaTeX reports when a long table splits.
\makeatletter
\@ifpackagelater{longtable}{2026/05/31}{}{%
  \ExplSyntaxOn
  \def\LT@output{%
    \ifnum\outputpenalty <-\@Mi
      \ifnum\outputpenalty > -\LT@end@pen
        \LT@err{floats~ and~ marginpars~ not~ allowed~ in~ a~ longtable}\@ehc
      \else
        \setbox\z@\vbox{\unvbox\@cclv}%
        \ifdim \ht\LT@lastfoot>\ht\LT@foot
          \dimen@\pagegoal
          \advance\dimen@\ht\LT@foot
          \advance\dimen@-\ht\LT@lastfoot
          \ifdim\dimen@<\ht\z@
            \setbox\@cclv\vbox{%
              \unvbox\z@\copy\LT@foot
              \vskip 0pt plus \maxdimen minus \normalbaselineskip
            }%
            \@makecol
            \@expl@@@mark@update@singlecol@structures@@
            \@outputpage
            \global\vsize\@colroom
            \setbox\z@\vbox{\box\LT@head}%
          \fi
        \fi
        \unvbox\z@\box\ifvoid\LT@lastfoot\LT@foot\else\LT@lastfoot\fi
        \UseTaggingSocket{tbl/longtable/foot}%
      \fi
    \else
      \setbox\@cclv\vbox{%
        \unvbox\@cclv\copy\LT@foot
        \vskip 0pt plus \maxdimen minus \normalbaselineskip
      }%
      \UseTaggingSocket{tbl/longtable/foot}%
      \@makecol
      \@expl@@@mark@update@singlecol@structures@@
      \@outputpage
      \global\vsize\@colroom
      \copy\LT@head\nobreak
    \fi}%
  \ExplSyntaxOff
}
\makeatother

\usepackage{mdframed}
\usepackage{wrapfig}

% Keep normal, uncompressed character spacing and standard single spacing.
\linespread{1.0}\selectfont
\setlength{\parindent}{0pt}
\setlength{\parskip}{0.35em}
\raggedbottom%

\setlength{\tabcolsep}{2.2pt}


% pa glossary colour
\definecolor{paWhite}{HTML}{FFFFFF}

\definecolor{ink}{HTML}{16324A}
\definecolor{accent}{HTML}{087E8B}

\definecolor{paBlue}{HTML}{244A73}
\definecolor{paMidBlue}{HTML}{5287C0}
\definecolor{paLightBlue}{HTML}{EEF4F8}
\definecolor{paTermBlue}{HTML}{DCEAF3}
\definecolor{paBorder}{HTML}{9EB6C8}

\definecolor{paGreen}{HTML}{438961}
\definecolor{paMidGreen}{HTML}{5A966F}
\definecolor{paTermGreen}{HTML}{DCEFE3}
\definecolor{paLightGreen}{HTML}{EEF7F1}
\definecolor{paGreenBorder}{HTML}{A4C5B1}

\definecolor{paGray900}{HTML}{222222}
\definecolor{paGray700}{HTML}{4D5661}
\definecolor{paGray500}{HTML}{7B8490}
\definecolor{paGray300}{HTML}{C9D0D6}
\definecolor{paGray100}{HTML}{F3F5F7}

\definecolor{paOrange}{HTML}{D97706} 
\definecolor{paLightOrange}{HTML}{FDF4E7} 
\definecolor{paTermOrange}{HTML}{F6E2BF} 
\definecolor{paOrangeBorder}{HTML}{D8B070}

\definecolor{LightOrange}{HTML}{FFB347}
\definecolor{Orange}{HTML}{FFA500}
\definecolor{DarkOrange}{HTML}{FF8C00}
\definecolor{BurntOrange}{HTML}{CC5500}
\definecolor{paRed}{HTML}{C0392B}


\setbeamercolor{normal text}{fg=ink,bg=white}
\setbeamercolor{structure}{fg=accent}
\setbeamercolor{frametitle}{fg=white,bg=planblue}
\setbeamerfont{frametitle}{size=\Large,series=\bfseries}
\setbeamertemplate{navigation symbols}{}
\setbeamertemplate{footline}{\hspace{6mm}\scriptsize Plan4ARI / A2 Review\hfill Draft for discussion\hspace{5mm}\insertframenumber\hspace{6mm}\vspace{4mm}}
\setbeamersize{text margin left=8mm,text margin right=8mm}
\setbeamertemplate{itemize items}[circle]
\setlength{\parskip}{5pt}

% ---------- Useful column types ----------
\newcolumntype{Y}{>{\raggedright\arraybackslash}X}
\newcolumntype{C}[1]{>{\centering\arraybackslash}p{#1}}
\newcolumntype{L}[1]{>{\raggedright\arraybackslash}p{#1}}
\newcolumntype{M}[1]{>{\centering\arraybackslash\bfseries}m{#1}}
\newcolumntype{J}[1]{>{\arraybackslash}p{#1}}

\newcolumntype{s}{>{\hsize=.8\hsize\raggedright\arraybackslash}X}

% ---------- pa concept box ----------
\newtcolorbox{paConcept}[1]{
    colback=paLightBlue,
    colframe=paBlue,
    colbacktitle=paBlue,
    coltitle=white,
    fonttitle=\bfseries,
    title={#1},
    boxrule=0.5pt,
    arc=1mm,
    left=1.5mm,
    right=1.5mm,
    top=1.5mm,
    bottom=1.5mm,
    before skip=3pt,
    after skip=3pt,
    breakable,
    enhanced
}

% ---------- pa concept box ----------
\newtcolorbox{paBox}[1][]{
    colback=paLightBlue,
    colframe=paBlue,
    boxrule=0.5pt,
    arc=1mm,
    left=1.5mm,
    right=1.5mm,
    top=1.5mm,
    bottom=1.5mm,
    before skip=3pt,
    after skip=3pt,
    breakable,
    enhanced,
    #1
}



% ---------- pa concept box ----------
\newtcolorbox{paBoxGreen}[1][]{
    colback=paLightGreen,
    colframe=paGreen,
    boxrule=0.1pt,
    arc=0.5mm,
    left=1.5mm,
    right=1.5mm,
    top=1.5mm,
    bottom=1.5mm,
    before skip=6pt,
    after skip=6pt,
    breakable,
    enhanced,
    #1
}

% ---------- pa examnple box ----------
\newtcolorbox{paExampleBox}[1][]{
    fonttitle=\bfseries,
    fontupper=\itshape,
    colback=paLightBlue,
    colframe=paLightBlue,
    coltitle=paBlue,
    boxrule=0.1pt,
    arc=0.1mm,
    left=1.mm,
    right=1.1mm,
    top=1.mm,
    bottom=1.mm,
    before skip=3pt,
    after skip=3pt,
    breakable,
    enhanced,
    #1
}


% ---------- pa light box ----------
\newtcolorbox{paLightBox}[1][]{
    colback=paGray100,
    colframe=paGray100,
    boxrule=0.1pt,
    arc=0.1mm,
    left=1.mm,
    right=1.1mm,
    top=1.mm,
    bottom=1.mm,
    before skip=3pt,
    after skip=3pt,
    breakable,
    enhanced,
    #1
}

% ---------- pa light box ----------
\newtcolorbox{paLightBluesBox}[1][]{
    colback=paLightBlue,
    colframe=paLightBlue,
    boxrule=0.1pt,
    arc=0.1mm,
    left=1.mm,
    right=1.1mm,
    top=1.mm,
    bottom=1.mm,
    before skip=3pt,
    after skip=3pt,
    breakable,
    enhanced,
    #1
}


\newtcbox{\boxedmath}[1][]{%
    nobeforeafter, math upper, tcbox raise base,
    enhanced, boxrule=0.2pt,
    #1}


\renewcommand{\cellalign}{vh}
\renewcommand{\theadalign}{vh}
\input{styles/math.tex}
% chktex-file 44



%%%%%%%%%%%%%%%%%%%%%%%%%%%%%%%%%%%%%%%%%%%%%%%%%%%%%%%%%%%%
% Basic Style Macros
\newcommand{\paemph}[1]{%
    \textcolor{pagRed}{\textit{#1}}%
}

\newcommand{\patextbf}[1]{%
    \textcolor{pagBlue}{\textbf{#1}}%
}

\newcommand{\patextbfo}[1]{%
    \textcolor{pagOrange}{\textbf{#1}}%
}
\newcommand{\placeholder}[1]{\textit{[#1]}}

\newcommand{\blueitem}{%
\textcolor{pagBlue}{\rule{0.8ex}{0.8ex}}% 
}
%%%%%%%%%%%%%%%%%%%%%%%%%%%%%%%%%%%%%%%%%%%%%%%%%%%%%%%%%%%%



%%%%%%%%%%%%%%%%%%%%%%%%%%%%%%%%%%%%%%%%%%%%%%%%%%%%%%%%%%%%
% Table Formatting
\newcommand{\parowcolored}[1]{%
    \rowcolor{#1}
    [\dimexpr\tabcolsep+0.2pt\relax]
    [\dimexpr\tabcolsep+0.2pt\relax]
}

\newcommand{\parowcolor}{%
    \parowcolored{pagMidBlue}%
}

\newcommand{\paheaderstrut}{%
  \rule[\dimexpr-\dp\strutbox-2pt\relax]{0pt}%
       {\dimexpr\ht\strutbox+\dp\strutbox+4pt\relax}%
}

\NewExpandableDocumentCommand{\pahcell}{o O{|l|} m}{%
    \IfNoValueTF{#1}{%
        \leavevmode\color{white}\paheaderstrut\textbf{#3}
    }{%
        \multicolumn{#1}{#2}{%
            \leavevmode\color{white}\paheaderstrut\textbf{#3}%
        }%
    }%
}

\NewDocumentCommand{\patableformat}{O{1.0} O{2pt} O{2pt}}{%
    \small
    \centering
    \renewcommand{\arraystretch}{#1}
    \setlength{\tabcolsep}{#2}
    \setlength{\extrarowheight}{#3}
    \arrayrulecolor{pagBorder}
}


\NewDocumentCommand{\patableformatnormalsize}{O{1.0} O{2pt} O{2pt}}{%
    \normalsize
    \centering
    \renewcommand{\arraystretch}{#1}
    \setlength{\tabcolsep}{#2}
    \setlength{\extrarowheight}{#3}
    \arrayrulecolor{pagBorder}
}

\newcommand{\rowspace}{\rule{0pt}{2ex}}


\NewDocumentCommand{\pafirstlinecell}{O{t} m m}{%
    \multirow[#1]{#2}{=}{%
        \parbox[t]{\dimexpr\linewidth-2\tabcolsep\relax}{%
            \raggedright #3 
        }%
    }% 
}

%%%%%%%%%%%%%%%%%%%%%%%%%%%%%%%%%%%%%%%%%%%%%%%%%%%%%%%%%%%%



%%%%%%%%%%%%%%%%%%%%%%%%%%%%%%%%%%%%%%%%%%%%%%%%%%%%%%%%%%%%
% TOC Formatting
\newlength{\palabelwidth}

\NewDocumentCommand{\pasubsubsectiononeline}{O{pagMidGreen} m +g o}{%
  \par
  \begingroup
  \raggedright
  \settowidth{\palabelwidth}{\textbf{#2}}%
  \addtolength{\palabelwidth}{4pt}%
  \phantomsection
  \addcontentsline{toc}{subsubsection}{#2}%
  \nopagebreak[4]%
  \noindent
  \IfNoValueTF{#3}{%
    \begin{tabularx}{\linewidth}{%
      >{\raggedright\arraybackslash}X}%
      \parowcolored{#1}\pahcell{#2}\\%
    \end{tabularx}%
  }{%
    \begin{tabularx}{\linewidth}{%
      >{\raggedright\arraybackslash}p{\palabelwidth}
      >{\raggedright\arraybackslash}X}%
      \cellcolor{#1}\pahcell{#2}
      &%
      \leavevmode\color{#1}\paheaderstrut\textbf{#3}
      \IfValueT{#4}{%
        \hfill
        {\color{black}#4} 
        \kern4pt \vrule width 1.5pt
      }\\%
      \IfNoValueT{#4}{\arrayrulecolor{#1}\hline}
    \end{tabularx}%
    \arrayrulecolor{pagBorder}%
  }%
  \par
  \endgroup
}

\NewDocumentCommand{\pasubsubsectionboxed}{O{pagMidGreen} m +g}{%
  \par
  \begingroup
  \phantomsection
  \addcontentsline{toc}{subsubsection}{#2}%
  \nopagebreak[4]%
  \noindent
  \begin{tabular}{|>{\raggedright\arraybackslash}%
    p{\dimexpr\linewidth-2\tabcolsep-2\arrayrulewidth\relax}|}%
    \rowcolor{#1}%
    \leavevmode\color{white}\paheaderstrut\textbf{#2}\\%
    \IfNoValueF{#3}{%
      \leavevmode\color{#1}\paheaderstrut\textbf{#3}\\%
      \hline
    }%
  \end{tabular}%
  \arrayrulecolor{pagBorder}%
  \par
  \endgroup
}

\NewDocumentCommand{\pasubsubsectionnamedboxed}
  {O{pagMidBlue} O{pagMidBlue} m m m o}{%
    \par
    \begingroup
    \raggedright
    \settowidth{\palabelwidth}{\textbf{#3}}%
    \addtolength{\palabelwidth}{4pt}%
    \IfNoValueT{#6}{
        \phantomsection
        \addcontentsline{toc}{subsubsection}{#3: #4}%
        \nopagebreak[4]%
    }
    \noindent
    \arrayrulecolor{#1}
    \begin{tabularx}{\dimexpr\linewidth-2pt\relax}{%
        >{\raggedright\arraybackslash}p{\palabelwidth}
        >{\raggedright\arraybackslash}X}%
        \cellcolor{#1}\pahcell{#3} &
        \leavevmode\color{#1}\paheaderstrut\textbf{#4}%
        \par
        \hrule height \arrayrulewidth
        \kern1pt
        \color{#2}\noindent\paheaderstrut#5\\%
    \end{tabularx}%
    \arrayrulecolor{pagBorder}%
    \par
    \endgroup
}
\NewDocumentCommand{\pasubsubsection}{m o}{%
    \par
    \phantomsection
    \addcontentsline{toc}{subsubsection}{#1}%
    \nopagebreak[4]%
    \textcolor{pagBlue}{%
        \textbf{#1}%
        \IfValueT{#2}{: \textbf{#2}}%
        \IfNoValueT{#2}{.}% 
    }% 
}

\newcounter{paRQ}
\NewExpandableDocumentCommand{\paRQrow}{m m}{%
  \IfValueT{#2}{%
    \IfBlankF{#2}{%
      \stepcounter{paRQ}%
      \patextbf{RQ\arabic{paRQ}} &
      \leavevmode\textit{#2}\\%
    }%
  }%
}

\newlength{\padefaultlabelwidth}
\NewDocumentCommand{\pasubsubsectionobjective}
  {O{pagMidGreen} m m m o o o o o}{%
    \par
    \begingroup
    \raggedright
    \settowidth{\palabelwidth}{\textbf{#2}}%
    \settowidth{\padefaultlabelwidth}{\textbf{RQ10}}%
    \ifdim\palabelwidth<\padefaultlabelwidth%
        \setlength{\palabelwidth}{\padefaultlabelwidth}%
    \fi%
    \addtolength{\palabelwidth}{4pt}%
    \phantomsection
    \addcontentsline{toc}{subsubsection}{#2: #3}%
    \nopagebreak[4]%
    \noindent
    \arrayrulecolor{#1}%
    \setlength{\extrarowheight}{1pt}%
    \begin{tabularx}{\linewidth}{%
        >{\raggedright\arraybackslash}p{\palabelwidth}
        >{\raggedright\arraybackslash}X}%
    \cellcolor{#1}\pahcell{#2} & 
    \leavevmode\color{#1}\paheaderstrut\textbf{#3}%
        \par
        \hrule height \arrayrulewidth
        \kern1pt
        \noindent\paheaderstrut#4\\%
        \paRQrow{#1}{#5}%
        \paRQrow{#1}{#6}%
        \paRQrow{#1}{#7}%
        \paRQrow{#1}{#8}%
        \paRQrow{#1}{#9}%
    \end{tabularx}%
    \arrayrulecolor{pagBorder}%
    \par
    \endgroup
}

%%%%%%%%%%%%%%%%%%%%%%%%%%%%%%%%%%%%%%%%%%%%%%%%%%%%%%%%%%%%%%%%

%%%%%%%%%%%%%%%%%%%%%%%%%%%%%%%%%%%%%%%%%%%%%%%%%%%%%%%%%%%%




%%%%%%%%%%%%%%%%%%%%%%%%%%%%%%%%%%%%%%%%%%%%%%%%%%%%%%%%%%%%
% ---------- Project ID ----------
\NewDocumentCommand{\patextbfgraphic}{O{pagGreen} O{pagGreen} m m}{%
\IfBlankTF{#3}{%
  \begin{tikzpicture}[baseline=(txt.base)]
    \node[inner xsep=0pt, inner ysep=0pt, text=#2, font=\bfseries] (txt) {#4};
    \draw[#1, line width=0.8pt] (txt.south west) -- (txt.south east);
  \end{tikzpicture}%
}{%
  \begin{tikzpicture}[baseline=(txt.base)]
    \node[inner sep=2pt, fill=#1, text=white, font=\bfseries] (tag) {#3};
    \node[anchor=west, inner xsep=0pt, inner ysep=0pt, text=#2, font=\bfseries] (txt) at ([yshift=-0.4pt]tag.east) {#4};
    \draw[#1, line width=0.8pt] (tag.south west) -- (txt.south east);
  \end{tikzpicture}%
}%
}

\newcommand{\ProjectName}{SMARTHER}
\newcommand{\CallIdentifier}{[HORIZON-EIC-2026-PATHFINDERCHALLENGES-01--03]}
\newcommand{\CallName}{[DeepRAP:\@ Deep Reasoning, Abstraction \& Planning]}
% \newcommand{\ProposalTitle}{%
% \begin{center}
% \begin{minipage}{0.15\textwidth}
% \includegraphics[width=\textwidth]{images/logo.png}
% \end{minipage}
% \begin{minipage}{0.7\textwidth}
% \centering
%   %\makebox[\textwidth][c]{%
%     \patextbfgraphic[pagOrange][black]
%       {}
%       {\patextbfo{S}emantic \patextbfo{M}etacognitive \patextbfo{A}utonomy Ha\patextbfo{R}ness for \patextbfo{T}rustworthy,}%
%   %}
%   \\[6pt]
%   %\makebox[\textwidth][c]{%
%     \patextbfgraphic[pagOrange][black]
%       {}
%       {\patextbfo{H}uman-Aware, \patextbfo{E}xplainable \patextbfo{R}obotics}%
%   %}%
% \end{minipage}
% \end{center}
% }

\newcommand{\ProposalTitle}{%
\includegraphics[width=0.95\textwidth]{images/logo9.png}
}


\DeclareRobustCommand{\pag}{%
  \textcolor{pagOrange}{\textbf{\ProjectName}}%
}
%%%%%%%%%%%%%%%%%%%%%%%%%%%%%%%%%%%%%%%%%%%%%%%%%%%%%%%%%%%%



%%%%%%%%%%%%%%%%%%%%%%%%%%%%%%%%%%%%%%%%%%%%%%%%%%%%%%%%%%%%
% Work-package description block. Duplicate the call once per WP.
\newcommand{\WorkPackageDescription}[4]{%
  \par
  \begingroup
  \setlength{\LTpre}{0pt}%
  \setlength{\LTpost}{0pt}%
  %
  \phantomsection
  \addcontentsline{toc}{subsection}{Table 3.1a - #1: #2}%
  \nopagebreak[4]%
  
  \begin{xltabular}{\textwidth}{|
    >{\raggedright\arraybackslash\bfseries}m{0.26\textwidth}|
    X|}
    \hline%
    \cellcolor{pagMidBlue}\color{white}Work package number & #1 \\\hline%
    \cellcolor{pagMidBlue}\color{white}Work package title & #2 \\\hline%
  \end{xltabular}
  \vspace{4pt}
  \begin{xltabular}{\textwidth}{|J{0.99\textwidth}|}
    \hline%
    \patextbf{Objectives.}
    #3\\[1.5em]\hline%
  \end{xltabular}
  \vspace{4pt}\begin{tcolorbox}[
  breakable,
  colback=white,
  colframe=pagBlue,
  boxrule=0.4pt,
  sharp corners,
  boxsep=0pt,
  left=3pt,
  right=3pt,
  top=3pt,
  bottom=3pt,
  before skip=2pt,
  after skip=18pt
]
    \setlength{\emergencystretch}{1em}%
    \patextbf{Description of work.}
    #4
  \end{tcolorbox}

  \endgroup
}
%%%%%%%%%%%%%%%%%%%%%%%%%%%%%%%%%%%%%%%%%%%%%%%%%%%%%%%%%%%%%%%%


\renewcommand{\tabularxcolumn}[1]{>{\raggedright\arraybackslash}p{#1}}
\newcommand{\src}[1]{\vfill{\tiny\color{gray}Source: #1\par}}
\newcommand{\lead}[1]{{\large\color{accent}\textbf{#1}}\par\medskip}
\newcommand{\items}{\setlength{\itemsep}{7pt}}


% Display short DOI values while linking them to their resolver URLs.
\DeclareFieldFormat{doi}{%
  \mkbibacro{DOI}\addcolon\space
  \ifhyperref
    {\href{https://doi.org/#1}{\nolinkurl{#1}}}
    {\nolinkurl{#1}}%
}

\DeclareFieldFormat{year}{\textcolor{DarkOrange}{\bfseries#1}}
\DeclareFieldFormat{date}{\textcolor{DarkOrange}{\bfseries#1}}
\DeclareFieldFormat{labeldate}{\textcolor{DarkOrange}{\bfseries#1}}


\title[Plan4ARI / A2 Review]{Plan4ARI\\\large Review of Activity A2}
\subtitle{Low-code programming with generative and deliberative AI}
\author{Plan4ARI}
\institute{Methodological update and technical specification alignment}
\date{Draft for technical discussion / 23 September 2026}
\hypersetup{pdftitle={Plan4ARI~-~Review of Activity A2},pdfauthor={Plan4ARI},pdfsubject={Proposed revision of the technical specification}}

\begin{document}

\begin{frame}
  \titlepage%
\end{frame}

\begin{frame}{Scope of the review}
\lead{Retain the objectives and clarify the implementation approach}
\begin{itemize}\items%
\item Identify A2 commitments in the technical specification and Annex A.
\item Assess the changes proposed in Scientific Plan D2.1 v0.11 and the Research Proposal dated 9 July 2026.
\item Map each change to tasks, expected outputs and verification methods.
\item Identify open decisions and inconsistencies across the documents.
\end{itemize}
\medskip
\textbf{Evidence status:} the proposal documents describe planned work. The M12 mockup is confirmed as a planned output; experiments and journal submission are not presented as completed.
\src{Technical specification, Activity 2; D2.1 v0.11, ``Purpose and boundary''; Research Proposal, abstract.}
\end{frame}


% Budget analysis: first three of four project years; CNR = AI+RO.
\begin{frame}{Budget analysis: scope and assumptions}


\begin{table}[H]\centering\scriptsize\caption{Mesi-persona per WP -- intero progetto.}\label{tab:wp-w-tot}
\begin{tabular}{llrrrr}\toprule\textbf{WP} & \textbf{Soggetto} & \textbf{Fund. Research} & \textbf{Ind. Research} & \textbf{Exp. Development} & \textbf{Totale} \\\midrule
WP2 & HI & 3.81 & 62.40 & 0.00 & 66.21 \\
 & AI & 1.96 & 3.91 & 0.00 & 5.87 \\
 & RO & 11.05 & 23.38 & 0.00 & 34.43 \\
 & \textit{Totale WP} & \textit{16.81} & \textit{89.70} & \textit{0.00} & \textit{106.51} \\\midrule
WP3 & HI & 3.81 & 62.40 & 0.00 & 66.21 \\
 & AI & 1.96 & 3.91 & 0.00 & 5.87 \\
 & RO & 3.68 & 7.79 & 0.00 & 11.48 \\
 & \textit{Totale WP} & \textit{9.44} & \textit{74.11} & \textit{0.00} & \textit{83.56} \\\midrule
WP4 & HI & 3.81 & 62.40 & 0.00 & 66.21 \\
 & AI & 1.96 & 3.91 & 0.00 & 5.87 \\
 & RO & 9.21 & 19.49 & 0.00 & 28.69 \\
 & \textit{Totale WP} & \textit{14.97} & \textit{85.80} & \textit{0.00} & \textit{100.77} \\
\bottomrule\end{tabular}
\end{table}

\end{frame}

\begin{frame}{Budget analysis: scope and assumptions}


\begin{table}[H]\centering\scriptsize
\begin{tabular}{llrrrr}\toprule\textbf{WP} & \textbf{Soggetto} & \textbf{Fund. Research} & \textbf{Ind. Research} & \textbf{Exp. Development} & \textbf{Totale} \\\midrule
WP5 & HI & 3.81 & 62.40 & 0.00 & 66.21 \\
 & AI & 1.96 & 3.91 & 0.00 & 5.87 \\
 & RO & 9.21 & 19.49 & 0.00 & 28.69 \\
 & \textit{Totale WP} & \textit{14.97} & \textit{85.80} & \textit{0.00} & \textit{100.77} \\\midrule
WP6 & HI & 3.81 & 62.40 & 0.00 & 66.21 \\
 & AI & 1.96 & 3.91 & 0.00 & 5.87 \\
 & RO & 3.68 & 7.79 & 0.00 & 11.48 \\
 & \textit{Totale WP} & \textit{9.44} & \textit{74.11} & \textit{0.00} & \textit{83.56} \\\midrule
WP7 & HI & 0.00 & 0.00 & 79.91 & 79.91 \\
 & AI & 0.00 & 0.00 & 9.78 & 9.78 \\
 & RO & 0.00 & 0.00 & 6.89 & 6.89 \\
 & \textit{Totale WP} & \textit{0.00} & \textit{0.00} & \textit{96.58} & \textit{96.58} \\
\midrule\textbf{Totale} & & \textbf{65.63} & \textbf{409.53} & \textbf{96.58} & \textbf{571.74} \\
\bottomrule\end{tabular}
\par\smallskip\footnotesize Le righe ``Totale WP'' e ``Totale'' sono calcolate in questa trascrizione (non presenti nel file).
\end{table}

\end{frame}

\begin{frame}{Budget analysis: Y1}
  
\begin{table}[H]\centering\scriptsize\caption{Mesi-persona per WP -- anno 1.}\label{tab:wp-w-y1}
\begin{tabular}{llrrrr}\toprule\textbf{WP} & \textbf{Soggetto} & \textbf{Fund. Research} & \textbf{Ind. Research} & \textbf{Exp. Development} & \textbf{Totale} \\\midrule
WP2 & HI & 1.90 & 20.84 & 0.00 & 22.75 \\
 & AI & 0.98 & 0.78 & 0.00 & 1.76 \\
 & RO & 5.52 & 7.01 & 0.00 & 12.54 \\
 & \textit{Totale WP} & \textit{8.40} & \textit{28.64} & \textit{0.00} & \textit{37.04} \\\midrule
WP3 & HI & 1.90 & 20.84 & 0.00 & 22.75 \\
 & AI & 0.98 & 0.78 & 0.00 & 1.76 \\
 & RO & 1.84 & 2.34 & 0.00 & 4.18 \\
 & \textit{Totale WP} & \textit{4.72} & \textit{23.96} & \textit{0.00} & \textit{28.69} \\\midrule
WP4 & HI & 1.90 & 20.84 & 0.00 & 22.75 \\
 & AI & 0.98 & 0.78 & 0.00 & 1.76 \\
 & RO & 4.60 & 5.85 & 0.00 & 10.45 \\
 & \textit{Totale WP} & \textit{7.48} & \textit{27.47} & \textit{0.00} & \textit{34.95} \\
\bottomrule\end{tabular}
\end{table}
\end{frame}



\begin{frame}{Budget analysis: Y1}
  
\begin{table}[H]\centering\scriptsize
\begin{tabular}{llrrrr}\toprule\textbf{WP} & \textbf{Soggetto} & \textbf{Fund. Research} & \textbf{Ind. Research} & \textbf{Exp. Development} & \textbf{Totale} \\\midrule
WP5 & HI & 1.90 & 20.84 & 0.00 & 22.75 \\
 & AI & 0.98 & 0.78 & 0.00 & 1.76 \\
 & RO & 4.60 & 5.85 & 0.00 & 10.45 \\
 & \textit{Totale WP} & \textit{7.48} & \textit{27.47} & \textit{0.00} & \textit{34.95} \\\midrule
WP6 & HI & 1.90 & 20.84 & 0.00 & 22.75 \\
 & AI & 0.98 & 0.78 & 0.00 & 1.76 \\
 & RO & 1.84 & 2.34 & 0.00 & 4.18 \\
 & \textit{Totale WP} & \textit{4.72} & \textit{23.96} & \textit{0.00} & \textit{28.69} \\\midrule
WP7 & HI & 0.00 & 0.00 & 0.00 & 0.00 \\
 & AI & 0.00 & 0.00 & 0.00 & 0.00 \\
 & RO & 0.00 & 0.00 & 0.00 & 0.00 \\
 & \textit{Totale WP} & \textit{0.00} & \textit{0.00} & \textit{0.00} & \textit{0.00} \\
\midrule\textbf{Totale} & & \textbf{32.82} & \textbf{131.51} & \textbf{0.00} & \textbf{164.33} \\
\bottomrule\end{tabular}
\par\smallskip\footnotesize Le righe ``Totale WP'' e ``Totale'' sono calcolate in questa trascrizione (non presenti nel file).
\end{table}
\end{frame}



\begin{frame}{Baseline commitments for A2}
\begin{columns}[T]
\begin{column}{.48\textwidth}
\lead{Application objective}
Enable process experts to program robotic applications through natural language.\par
Integrate the modules into COMAU MI.RA/Dexter.\par
Activity A2: \textbf{M1--M36}.
\end{column}
\begin{column}{.48\textwidth}
\lead{Targets}
\begin{itemize}\items
\item \textbf{66\%} reduction in programming time.
\item Programming by workers without prior programming training.
\item A fleet of \textbf{3 mobile manipulators} or \textbf{4 industrial robotic arms}.
\end{itemize}
\end{column}
\end{columns}
\src{Annex A, Section A.3.4, p. 10, and Activity 2, p. 36. The baseline uses ``not pre-trained worker''.}
\end{frame}

\begin{frame}{Original approach and proposed revision}
\small\renewcommand{\arraystretch}{1.4}
\begin{tabularx}{\textwidth}{@{}p{.22\textwidth}XX@{}}
\toprule
Stage & Original approach & D2.1 v0.11 proposal\\\midrule
Knowledge & Knowledge base from process descriptions & Capability registry with evidence and verification\\
Formalisation & PDDL generation using a lightweight LLM & Typed intermediate representation and compilation\\
Planning & BTs and digital-twin optimisation & Solve and validate the plan before BT generation\\
Optimisation & MCTS search in simulation & Solver portfolio, scheduling and geometric checks\\
Execution & BT execution and online replanning & Skill grounding, effect verification and recovery\\\bottomrule
\end{tabularx}
\src{Annex A, Section A.3 and Tasks 2.2--2.5; D2.1 v0.11, Sections 1.1, 6--8. Changes are proposed.}
\end{frame}

\begin{frame}{Scientific contribution to make explicit}
\lead{Symbolic plans must correspond to executable capabilities}
\begin{itemize}\items
\item An available ROS2 interface does not, by itself, specify a skill's preconditions, effects or constraints.
\item The capability registry combines runtime interfaces, documentation, MI.RA/Dexter blocks and verification evidence.
\item Grounding binds each planned action to a concrete skill with consistent parameters and resources.
\item Runtime monitoring compares expected effects with observed state and triggers recovery when required.
\end{itemize}
\src{D2.1 v0.11, Sections 1 and 7; Research Proposal, Sections 2.1--2.2. Summary of the proposed research.}
\end{frame}

\begin{frame}{Three research directions in the A2 proposal}
\small\renewcommand{\arraystretch}{1.35}
\begin{tabularx}{\textwidth}{@{}p{.26\textwidth}XX@{}}
\toprule
Direction & Research question & Proposed task mapping\\\midrule
1. Available domain & How can a request become a valid plan grounded in ROS2 skills? & T2.2, T2.3 and T2.5\\
2. Skill discovery & How can a domain be built or adapted from the system and its evidence sources? & T2.2, supported by T2.6\\
3. Scene graphs & How can relations between parts support disassembly planning and state verification? & T2.3--T2.5, with A3/A4 interfaces\\\bottomrule
\end{tabularx}
\par\vspace{3mm}
Direct BT generation from scene graphs is an exploratory research direction. Its integration into the validated planning pipeline requires explicit definition.
\src{Research Proposal, Sections 2.1--2.3; D2.1 v0.11, Sections 7 and 10. Task mapping proposed for this review.}
\end{frame}

\begin{frame}{Proposed functional architecture}
\small\renewcommand{\arraystretch}{1.3}
\begin{tabularx}{\textwidth}{@{}p{.25\textwidth}XX@{}}
\toprule
Stage & Output & Required checks\\\midrule
Operational knowledge & Skill registry and scene state & Provenance, availability and consistency\\
Interpretation & Typed goals and constraints & Ambiguity, objects and types\\
Formalisation & PDDL domain and problem & Model syntax and semantics\\
Planning & Plan and resource allocation & Constraints, validity and solution quality\\
Execution preparation & BT and executor bindings & Parameters, resources and twin feasibility\\
Execution & Execution trace and updated state & Preconditions, effects and timing\\\bottomrule
\end{tabularx}
\src{D2.1 v0.11, Sections 7--8. Proposed functional sequence, not an implemented architecture.}
\end{frame}

\begin{frame}{Proposed revision of T2.2--T2.4}
\small\renewcommand{\arraystretch}{1.35}
\begin{tabularx}{\textwidth}{@{}p{.16\textwidth}XX@{}}
\toprule
Task and period & Proposed change & Required evidence\\\midrule
T2.2\newline M13--M24 & Capability discovery, lightweight models, typed representation and PDDL compilation & Verified registry, reference problems and comparison with direct generation\\
T2.3\newline M18--M24 & BTs derived from validated plans. Explicit temporal handling using STN/STNU & Plan-to-BT correspondence and variable-duration tests\\
T2.4\newline M22--M30 & Solver selection by problem type. Twin-based feasibility and cost estimation & Quality against a baseline, motion feasibility and computational cost\\\bottomrule
\end{tabularx}
\src{Annex A, p. 37; D2.1 v0.11, Sections 8 and 10. Baseline task periods retained in this proposal.}
\end{frame}

\begin{frame}{Proposed revision of T2.5--T2.8}
\small\renewcommand{\arraystretch}{1.3}
\begin{tabularx}{\textwidth}{@{}p{.17\textwidth}XX@{}}
\toprule
Task and period & Proposed change & Required evidence\\\midrule
T2.5\newline M28--M36 & State and effect checks, bounded plan repair and replanning & Perturbation tests, recovery times and required interventions\\
T2.6\newline M28--M36 & MI.RA/Dexter block binding and module integration & Integrated workflows and interface tests\\
T2.7\newline M13--M36 & Versioning of software, models, prompts and benchmarks & Reproducible tests and release traceability\\
T2.8\newline M7--M12 & Initial benchmark definition, with proposed ongoing evaluation & Protocol, test cases and end-user evaluation\\\bottomrule
\end{tabularx}
\src{Annex A, pp. 37--38; D2.1 v0.11, Section 10. Extending T2.8 requires explicit schedule alignment.}
\end{frame}

\begin{frame}{Deliverables and schedule dependencies}
\small\renewcommand{\arraystretch}{1.35}
\begin{tabularx}{\textwidth}{@{}p{.13\textwidth}XX@{}}
\toprule
Deadline & Baseline commitment & Issue to resolve\\\midrule
D2.1 / M12 & Scientific plan and first running mockup of the low-code module & Confirmed for M12. D2.1 v0.11 includes the mockup described in the Research Proposal and its demonstration evidence\\
D2.2 / M24 & First time-variance-aware planner running in MI.RA/Dexter & Minimum integration is needed by M24. T2.6 starts at M28; the v0.9 roadmap places integration at M25--M30\\
D2.3 / M36 & Optimal task and motion planning with online replanning & Link the demonstrator to KPI evidence and acceptance criteria\\\bottomrule
\end{tabularx}
\src{Annex A, pp. 38 and 46; D2.1 v0.11, ``Purpose and boundary'' and Section 10. Issues identified through comparison.}
\end{frame}

\begin{frame}{First running mockup at M12}
\lead{D2.1 includes the scientific plan and the running mockup}
\begin{itemize}\items
\item Operator GUI, natural-language request, optional context and an available PDDL domain.
\item JSON problem generation, deterministic PDDL compilation, validation and repair.
\item Planner portfolio and action grounding through ROS2 introspection and PlanSys2 executors.
\item Initial single-robot pick-and-place demonstration with versioned artefacts and execution evidence.
\end{itemize}
\small Skill-registry generation and mechanical scene graphs extend the research programme. The M12 report will state which prototype functions are implemented and tested.
\src{Research Proposal, Section 2.1, Figure 1 and Section 3; D2.1 v0.11, Section 6.11. M12 availability confirmed in the project update.}
\end{frame}

\begin{frame}{Scientific manuscript submission plan}
\lead{Submission target: 23 October 2026}
\begin{itemize}\items
\item Paper on the mockup, neuro-symbolic architecture and initial experimental evaluation.
\item Candidate journal: IEEE Transactions on Automation Science and Engineering (IEEE T-ASE) or IEEE Transactions on Industrial Informatics (IEEE TII).
\item Complete experiments, reproducibility materials, failure analysis and co-author review before submission to one selected journal.
\item Record submission acknowledgement separately from any subsequent acceptance or publication.
\end{itemize}
\small The calendar target is within one month of the 23 September update. It does not redefine the project M12 date.
\src{Project planning update, 23 September 2026; D2.1 v0.11, Sections 8.7 and 10.5. Planned submission.}
\end{frame}

\begin{frame}{Document corrections for discussion}
\small\renewcommand{\arraystretch}{1.35}
\begin{tabularx}{\textwidth}{@{}p{.22\textwidth}XX@{}}
\toprule
Issue & Evidence & Proposed correction\\\midrule
T2.1 period & Original specification: M7--M12. Revised A2 table: M1--M12 & Aligned to M1--M12 in D2.1 v0.11, following the proposed specification revision\\
T2.4 input & Text assigns BT output to T2.2. BT generation belongs to T2.3 & Correct the reference and specify the actual input\\
M24/M36 milestones & The milestone list still cites D2.1 & Check references to D2.2 and D2.3. Also clarify the M48 reference\\
T2.8 period & M7--M12, although the text includes platform evaluation & Separate initial benchmark design from later evaluation campaigns\\\bottomrule
\end{tabularx}
\src{Annex A, pp. 37--38 and 46; technical specification, Activity 2 and milestones; D2.1 v0.11, Sections 2.3 and 10.}
\end{frame}

\begin{frame}{Baseline KPIs and proposed metrics}
\begin{columns}[T]
\begin{column}{.48\textwidth}
\lead{Baseline commitments}
\begin{itemize}\items
\item Programming time reduction: 66\%.
\item Workers without prior training.
\item 3 mobile manipulators or 4 arms.
\item Knowledge base: 10 contexts.
\item PDDL, BT and task planning: ``+100 contexts'' for each category.
\end{itemize}
\end{column}
\begin{column}{.48\textwidth}
\lead{Additional D2.1 metrics}
\begin{itemize}\items
\item Capability discovery accuracy.
\item Model and plan validity.
\item Actions grounded in executable skills.
\item Planning quality and latency.
\item Robustness and human interventions.
\end{itemize}
\end{column}
\end{columns}
\medskip\small New thresholds in v0.9 are proposals. The definition of a ``context'' and the counting protocol require agreement.
\src{Annex A, Section A.3.4, p. 10; D2.1 v0.11, Section 9 and metric table.}
\end{frame}

\begin{frame}{Proposed evaluation protocol}
\lead{Measure the full process up to an executable workflow}
\begin{enumerate}\items
\item Define cases, initial conditions and correct outcomes with domain experts.
\item Compare direct generation, a validator-based pipeline and a capability-registry-based pipeline.
\item Measure initial cell configuration, task programming and human correction effort separately.
\item Test duration changes, missing objects and unavailable skills.
\item Evaluate the integrated workflow with operators at the target robot fleet scale.
\end{enumerate}
\small
Programming time reduction: $100\,(T_{\mathrm{baseline}}-T_{\mathrm{A2}})/T_{\mathrm{baseline}}$.\par
Use identical tasks and completion criteria in both conditions.
\src{Review proposal based on Annex A, Section A.3.4 and T2.8; D2.1 v0.11, Section 9; Research Proposal, Section 3.}
\end{frame}

\begin{frame}{Interfaces and responsibilities to agree}
\small\renewcommand{\arraystretch}{1.2}
\begin{tabularx}{\textwidth}{@{}>{\raggedright\arraybackslash}p{.25\textwidth}X@{}}
\toprule
Interface & Required agreement\\\midrule
CNR / planning specialists & Capability models, formalisation, solvers, validation and prototypes\\
COMAU / MI.RA/Dexter & Block catalogue, interfaces, industrial requirements and integration deadlines\\
A3 / A4 & Perceived state, motion feasibility checks and execution feedback\\
A6 / users & Comprehensibility, workload, interventions and task completion\\
A2 lead / PI & Methodological choices, ownership, acceptance criteria and change management\\\bottomrule
\end{tabularx}
\par\vspace{3mm}
Proposed boundary: cognitive functions in the Open IPC; hard-real-time control and safety-critical functions in the Core IPC.
\src{D2.1 v0.11, Sections 2.4 and 10, including the roles table. Proposed assignments require confirmation.}
\end{frame}

\begin{frame}{Potential COMAU design contributions (1/3)}
\small
\textbf{Evidence status:} contributions to confirm and document, not verified completed activities.
\medskip
\renewcommand{\arraystretch}{1.35}
\begin{tabularx}{\textwidth}{@{}>{\raggedright\arraybackslash}p{.23\textwidth}XX@{}}
\toprule
Potential COMAU contribution & Effect on the method and mockup & Evidence to retrieve\\\midrule
Dexter programming workflow analysis & Identify decisions that currently require specialist expertise and those suitable for automation. & Annotated demonstrations, example workflows, meeting notes and emails.\\
Block-library analysis & Specify how a block becomes a plannable skill: parameters, preconditions, effects, resources and errors. & Module descriptions, configuration examples and technical annotations.\\
Industrial requirements definition & Inform structured generation, validation and explicit handling of missing capabilities. & Requirements, architecture review comments and documented decisions.\\
\bottomrule
\end{tabularx}
\src{Contribution framework based on the Host Institution role and the project team's description of Dexter. D2.1, Section 9.}
\end{frame}

\begin{frame}{Potential COMAU design contributions (2/3)}
\small
\textbf{Evidence status:} contributions to confirm and document, not verified completed activities.
\medskip
\renewcommand{\arraystretch}{1.35}
\begin{tabularx}{\textwidth}{@{}>{\raggedright\arraybackslash}p{.23\textwidth}XX@{}}
\toprule
Potential COMAU contribution & Effect on the method and mockup & Evidence to retrieve\\\midrule
Use-case definition & Define realistic process sequences, constraints, variants and mockup failure conditions. & Use-case specification, process examples and test scenarios.\\
Future interface design & Specify the mapping between symbolic actions, ROS2 skills and Dexter modules for subsequent integration. & Interface diagrams, preliminary mappings and compatibility assessments.\\
Mockup review and acceptance criteria & Assess industrial relevance and define how the usefulness of the prototype will be measured. & Demo feedback, annotated tests, issue lists and resulting corrections.\\
\bottomrule
\end{tabularx}
\src{Contribution framework based on the Host Institution role and the project team's description of Dexter. D2.1, Section 9.}
\end{frame}


\begin{frame}{Potential COMAU design contributions (3/3)}
\lead{In other words}
COMAU should
\begin{itemize}
\item Define the Knowledge base and skill registry for the low-code programming framework.
\item Define the operational capability models and their verification evidence.
\item Define the use cases for the low-code programming framework.
\item Design the interface for the framework.
\item Review and define the acceptance criteria for the mockup.
\end{itemize}
\url{
https://gist.github.com/karpathy/442a6bf555914893e9891c11519de94f
}
\par
\end{frame}


\begin{frame}{Proposed wording for the A2 objective}
\lead{Working text for the revised technical specification}
Activity A2 will develop a low-code programming framework combining generative AI and deliberative planning, intended for process experts and integrated into MI.RA/Dexter.\par
The framework will translate natural-language requests into verifiable planning models, using an operational capability registry and a scene-state representation. Planned actions will be bound to executable skills and checked against temporal constraints, resource availability and motion feasibility.\par
Execution will include effect monitoring and replanning. Evaluation will document the baseline targets for programming time reduction, operator accessibility and multi-robot scale.
\src{Proposed wording based on the technical specification, Activity 2, and D2.1 v0.11. This is not approved text.}
\end{frame}

\begin{frame}{Decisions required for the revision}
\begin{enumerate}\items
\item Confirm the pipeline based on a capability registry, validated models and executable skill bindings.
\item Define the role of exploratory work: disassembly scene graphs, MCTS and adapter generation.
\item Confirm the M12 mockup evidence package and the manuscript submission plan.
\item Align minimum integration with D2.2/M24 and clarify continued T2.8 evaluation.
\item Agree on benchmarks, ownership and acceptance criteria, keeping baseline KPIs distinct from additional metrics.
\end{enumerate}
\src{Summary of open issues identified through the document comparison.}
\end{frame}

\begin{frame}{Sources and document status}
\small
\begin{itemize}\setlength{\itemsep}{2pt}
\item \textbf{Annex A, FISA-2024-00099}: technical reference for objectives, KPIs, tasks and deliverables. Main passages: p. 10, pp. 36--38 and p. 46.
\item \textbf{Plan4ARI v3, final technical specification}: editable Activity 2 text and deliverable list used for comparison.
\item \textbf{D2.1 Updated Scientific Plan, v0.9}: revised internal draft including the M12 mockup, methodological updates and manuscript submission plan.
\item \textbf{A2 Research Project Proposal, 9 July 2026}: research proposal covering neuro-symbolic planning, skill discovery and scene graphs.
\item \textbf{budget\_COMAU\_v6.xlsx, Budget Per Year}: annual partner costs and funding; hypothetical effort at EUR 6,100/PM for COMAU and EUR 5,200/PM for CNR.
\item \textbf{Annexes B, C and D}: financial framework, project terms and identifiers. No budget amendment is proposed.
\end{itemize}
\medskip
This presentation compares the supplied documents. It does not include a new external verification of the literature cited in D2.1.
\end{frame}
\end{document}
